\begin{titlepage}


    \setlength{\centeroffset}{-0.5\oddsidemargin}
    \addtolength{\centeroffset}{0.5\evensidemargin}
    \thispagestyle{empty}

    \begin{center}
        \begin{minipage}{\textwidth}

            \centering
            % \includegraphics[width=0.7\textwidth]{imagenes/logo_ugr.jpg}\\[1.4cm]

            \textsc{ \Large PROYECTO FIN DE CARRERA\\[0.2cm]}
            \textsc{ INGENIERÍA EN INFORMÁTICA}\\[1cm]
            % Upper part of the page
            % 

            %si el proyecto tiene logo poner aquí
            \includegraphics[width=0.5\linewidth]{logo/white_logo_text_stroke.png}
            \vspace{0.5cm}

            % Title

            {\Huge\bfseries Gestión de Deuda Técnica y Diseño de Interfaz en Framework MLOps\\
            }
            \noindent\rule[-1ex]{\textwidth}{3pt}\\[3.5ex]
        \end{minipage}
    \end{center}

    \vspace{1.65cm}
    \begin{center}
        \begin{minipage}{\textwidth}
            \centering

            \textbf{Autor}\\ {Yeray López Ramírez (alumno)}\\[2.5ex]
            \textbf{Directores}\\
            {Alberto Guillén Perales (tutor)}\\[2cm]
            %\includegraphics[width=0.15\textwidth]{imagenes/tstc.png}\\[0.1cm]
            %\textsc{Departamento de Teoría de la Señal, Telemática y Comunicaciones}\\
            %\textsc{---}\\
            %Granada, mes de 201
        \end{minipage}
    \end{center}
    %\addtolength{\textwidth}{\centeroffset}
    \vspace{\stretch{2}}


\end{titlepage}




\cleardoublepage
\thispagestyle{empty}

\begin{center}
       {\large\bfseries Gestión de Deuda Técnica y Diseño de Interfaz en Framework MLOps (SnapperML)}\\
\end{center}
\begin{center}
       Yeray López Ramírez (alumno)\\
\end{center}

%\vspace{0.7cm}
\noindent{\textbf{Palabras clave}: Deuda técnica, MLOps, Interfaz Web, Artectos, Reproducibilidad, Framework, Bases de Datos, Experimentos, Optimización de Hiperparámetros, Configuración de experimento.}

\vspace{1.0cm}

\noindent{\textbf{Resumen}}

\vspace{0.5cm}

Este Trabajo Fin de Grado se centra en la resolución de deuda técnica y el desarrollo de una interfaz web para el framework MLOps SnapperML. El proyecto aborda la deuda técnica existente para mejorar la mantenibilidad, escalabilidad y rendimiento general de SnapperML. La metodología incluye la identificación y refactorización de áreas con alta complejidad acumulada, la actualización de dependencias y la implementación de pruebas automatizadas para garantizar la estabilidad del sistema. Además, se desarrolló una interfaz web utilizando tecnologías modernas de frontend que facilita la interacción del usuario, permitiendo una gestión fluida de las configuraciones y la ejecución de experimentos de aprendizaje automático.

\vspace{2.5mm}

SnapperML integra principios clave de MLOps, asegurando la reproducibilidad y una gestión eficiente de artefactos. La implementación abarca el framework basado en CLI, una API en Flask y un servidor web con Vite, junto con bases de datos contenedorizadas para almacenar los resultados de los experimentos. El proyecto también destaca estrategias para prevenir la acumulación futura de deuda técnica, siguiendo las mejores prácticas de ingeniería de software. En conjunto, las mejoras realizadas buscan proporcionar una herramienta robusta y accesible para la gestión de flujos de trabajo de aprendizaje automático, garantizando la sostenibilidad y usabilidad a largo plazo.
\cleardoublepage


\thispagestyle{empty}


\begin{center}
       {\large\bfseries Technical Debt Management and Web Design in MLOps Framework (SnapperML)}\\
\end{center}
\begin{center}
       Yeray López Ramírez (student)\\
\end{center}

%\vspace{0.7cm}
\noindent{\textbf{Keywords}: Technical Debt, MLOps, Web Interface, Artifacts, Reproducibility, Framework, Databases, Experiments, Hyperparameter Optimization, Experiment Configuration}

\vspace{1.2cm}

\noindent{\textbf{Abstract}}

\vspace{0.5cm}

This Final Degree Project focuses on the resolution of technical debt and the development of a user-friendly web interface for the SnapperML MLOps framework. The project addresses existing technical debt to enhance the maintainability, scalability, and overall performance of SnapperML. The approach includes identifying and refactoring areas of accumulated complexity, updating dependencies, and implementing a series of automated tests to ensure stability. Additionally, a web interface was developed using modern frontend technologies to facilitate user interaction, allowing for seamless configuration management and execution of machine learning experiments.

\vspace{2.5mm}

SnapperML integrates key MLOps principles, ensuring reproducibility and efficient artifact management. The implementation covers the CLI-based framework, a Flask API, and a Vite-based web server, coupled with containerized databases to store experimental results. The project also emphasizes strategies to prevent future technical debt, adhering to best software engineering practices. Overall, the improvements aim to provide a robust and accessible tool for managing machine learning workflows, ensuring long-term project sustainability and usability.

\chapter*{}
\thispagestyle{empty}

\noindent\rule[-1ex]{\textwidth}{2pt}

\vspace{1.3cm}

Yo, \textbf{Yeray López Ramírez}, alumno de la titulación Ingeniería Informática de la \textbf{Escuela Técnica Superior
       de Ingenierías Informática y de Telecomunicación de la Universidad de Granada}, con DNI 26050768W, autorizo la
ubicación de la siguiente copia de mi Trabajo Fin de Grado en la biblioteca del centro para que pueda ser
consultada por las personas que lo deseen.

\vspace{6cm}

\noindent Fdo: Yeray López Ramírez

\vspace{2cm}

\begin{flushright}
       Granada a 4 de noviembre de 2024 .
\end{flushright}


\chapter*{}
\thispagestyle{empty}

\noindent\rule[-1ex]{\textwidth}{2pt}

\vspace{1.2cm}

D. \textbf{Alberto Guillén Perales (tutor)}, Profesor del Departamento de Ingeniería de Computadores, Automática y Robótica de la Universidad de Granada.

\vspace{0.5cm}

\textbf{Informa:}

\vspace{0.5cm}

Que el presente trabajo, titulado \textit{\textbf{Gestión de Deuda Técnica y Diseño de Interfaz en Framework MLOps (SnapperML)}},
ha sido realizado bajo su supervisión por \textbf{Yeray López Ramírez (alumno)}, y autorizamos la defensa de dicho trabajo ante el tribunal
que corresponda.

\vspace{0.5cm}

Y para que conste, expiden y firman el presente informe en Granada a 4 de noviembre de 2024.

\vspace{1cm}

\textbf{Los directores:}

\vspace{5cm}

\noindent \textbf{Alberto Guillén Perales (tutor)}

\chapter*{Agradecimientos}
\thispagestyle{empty}

\vspace{1cm}

Quiero expresar mi más sincero agradecimiento a mi padre por su apoyo incondicional durante los estudios y a mi hermano por su ayuda que hizo que mi inicio en la universidad fuera más llevadero.

\vspace{2.5mm}

A la DEIIT por brindarme la oportunidad de prestar servicios a los estudiantes de la facultad. Haber podido contribuir al bienestar de mis compañeros hizo que mi paso por la universidad fuera aún más enriquecedor.

\vspace{2.5mm}

A la Escuela y a mis compañeros de clase, por retarme a ser la mejor versión de mí mismo, promoviendo un entorno de aprendizaje y colaboración que me impulsó a superarme constantemente.

\vspace{2.5mm}

A mis compañeros de piso, con quienes compartí gran parte de mi estancia. Gracias por las incontables horas de risas y compañía en el comedor y la cafetería, que llenaron nuestra convivencia de buenos momentos y recuerdos inolvidables.

\vspace{2.5mm}

A mis amigos más cercanos, quienes con su apoyo constante y su amistad convirtieron esta etapa universitaria en una experiencia inolvidable para mi, mucho más allá de lo que jamás podría haber imaginado. No encuentro palabras suficientes para expresar la gratitud que siento por haber contado con vosotros.

\vspace{2.5mm}

A mi novio por acompañarme tanto en los malos como en los buenos momentos y por mantenerme en el buen camino. Te admiro por tu determinación y esfuerzo por equilibrar nuestras metas personales y nuestra relación hasta el día de hoy. Sin duda, conocerte me cambió la vida y espero que la tuya también, te quiero. 

\vspace{2.5mm}

Finalmente, a mi tutor, a quien admiro profundamente por su verdadera pasión por la docencia y su dedicación a la formación de auténticos ingenieros. Agradezco de corazón tu orientación y compromiso, que marcaron una diferencia significativa en mi desarrollo académico y profesional.

\vspace{2.5mm}

A todos ellos, los recordaré siempre con cariño y gratitud, por haber sido parte esencial de esta etapa tan significativa de mi vida.